% OPTION 3 - a replacement for \subsection{Alternate System} (sec:Alternate) that keeps
% its purpose (easy optimization, gamma < 1, first-order weights, yet T(0) != 1) AND
% doubles as the anti-windup example, so that sec:antiwindup can refer to it and the
% drivetrain figure YH-AW-18 is not needed. Same plant as before; the loop crossover is
% moved from 1 rad/s to 0.3 rad/s, below the plant pole.
%
% Figures and numbers: ../make_YH-AW-3rd.py (or ../antiwindup_variants.m section D),
% re-synthesized with the repo's H-inf solver; Wp = 1e3 at DC crossing 0.3,
% Wd = 0.1 -> 20 crossing 0.3, Wu = 0.1 -> 30 crossing 30. If sec:antiwindup keeps its
% own drivetrain figure, drop the last paragraph here.
%
% Framing note (2026-09-16): this is the one case where the integrator-only scheme is
% sufficient. It is presented as a property of the realization (the explicit integrator
% state), not as Youla-H out-performing H-inf under anti-windup - with full-state
% conditioning the two controllers recover identically on every plant tested.

\subsection{Alternate System} \label{sec:Alternate}
An alternate system to demonstrate the value of the hybrid Youla-H gain tuning method
is shown below in equation \ref{eq:Gp_simple}.
\begin{equation} \label{eq:Gp_simple}
G_{P} = \frac{1}{s+1}
\end{equation}
The weighting functions chosen for this system are all first-order, compared to the
second-order slopes used in the drivetrain system, and place the crossover at
$0.3$\,rad/s, below the plant pole. The weights are relaxed, allowing for easy
optimization, but not so relaxed as to be unrealistic for a controller design
problem. Figure \ref{fig:TSY-simple} shows the weighting functions and the resulting
closed-loop transfer functions.
\begin{figure}[H]
    \includegraphics[width=1\linewidth]{antiwindup_study/figures/YH-TSY-3rd.pdf}
    \caption{Weighting functions and closed-loop transfer functions for the Youla-H
    controller on the simple plant.}
    \label{fig:TSY-simple}
\end{figure}
Here $\gamma = 0.987$. Despite the optimization succeeding, $T_H(0) = 0.99898$, a
deficiency of $0.1\%$ that a long ramp would integrate into a noticeable bias. Even
with a simple plant and relaxed weights, the $H_\infty$ optimization in the best-case
non-trivial example still leaves a minor tracking deficiency, and the Youla-H method
fixes it with no change to the weights; $M_2$ is unchanged at $0.9625$. See Appendix
A for example code.

With the crossover below the plant pole the controller is close to a PI: the split
of Equation \ref{eq:split} gives $k_I = 0.293$ with a remainder $R(s)$ whose DC gain
is only $0.23$ and whose poles lie at or beyond the plant pole, so the integrator
alone holds the low-frequency memory. This is the case in which the explicit
integrator state of the Youla-H realization is worth the most. Figure
\ref{fig:YH-AW-3rd} holds the reference at $3$, beyond an actuator limit
$|u|\le 1.5$, for $2.5$\,s and then drops it to $1$. With a clamp alone, $G_{CH}$
and $G_{CYH}$ alike remain on the rail for $0.7$\,s after the step-down, overshoot
to $1.43$, and take $12$\,s to settle; the $H_\infty$ controller winds up exactly as
the Youla-H one does, but offers no integrator state to condition. Back-calculating
the clamp excess into the Youla-H integrator releases the rail immediately and
settles in $2.5$\,s with no undershoot. No conditioning of the remaining states is
required here. Where the remainder does carry the low-frequency gain, as for the
drivetrain controller of Section \ref{sec:antiwindup}, full-state conditioning is
needed instead, and it then serves $G_{CH}$ and $G_{CYH}$ equally well; the
advantage of the Youla-H realization is the choice it makes available, together
with the exact DC tracking that survives saturation.
\begin{figure}[H]
    \includegraphics[width=1\linewidth]{antiwindup_study/figures/YH-AW-3rd.pdf}
    \caption{Saturated response of the simple plant (reference in light gray): clamp
    only, and back-calculation on the integrator state of the realization of
    Equation \ref{eq:split}.}
    \label{fig:YH-AW-3rd}
\end{figure}
